\documentclass[tikz,border=6pt]{standalone}
\usepackage{ctex}
\usepackage{amsmath}
\usepackage{pgfplots}
\pgfplotsset{compat=1.18}
\usetikzlibrary{calc,angles,quotes,arrows.meta,positioning,decorations.markings}
\begin{document}
\begin{tikzpicture}[scale=1.0]
  % 一般（锐角）三角形 ABC，边 a,b,c 分别对角 A,B,C
  \coordinate (A) at (0,0);
  \coordinate (B) at (6,0);
  \coordinate (C) at (4.2,3.4);

  % 三边
  \draw[very thick] (A) -- (B) -- (C) -- cycle;

  % 顶点的角标记
  \pic[draw,thick,"$A$",angle radius=0.85cm,angle eccentricity=1.55]
       {angle = B--A--C};
  \pic[draw,thick,"$B$",angle radius=0.85cm,angle eccentricity=1.55]
       {angle = C--B--A};
  \pic[draw,thick,"$C$",angle radius=0.85cm,angle eccentricity=1.7]
       {angle = A--C--B};

  % 顶点标注
  \fill (A) circle (1.8pt) node[below left] {$A$};
  \fill (B) circle (1.8pt) node[below right] {$B$};
  \fill (C) circle (1.8pt) node[above] {$C$};

  % 边标注：a 对 A（即 BC），b 对 B（即 CA），c 对 C（即 AB）
  \node[fill=white,inner sep=1.5pt] at ($(B)!0.5!(C)+(0.45,0.12)$) {$a$};
  \node[fill=white,inner sep=1.5pt] at ($(C)!0.5!(A)+(-0.42,0.1)$) {$b$};
  \node[fill=white,inner sep=1.5pt] at ($(A)!0.5!(B)+(0,-0.32)$) {$c$};

  % 半透明对应关系说明框
  \node[fill=cyan!14,fill opacity=0.78,text opacity=1,rounded corners,
        inner sep=4pt,align=left] at (0.05,2.85)
       {边对角：\\ $a$ 对 $A$\\ $b$ 对 $B$\\ $c$ 对 $C$};
\end{tikzpicture}
\end{document}
